
\subsubsection{6.1 Synthetic vs Real Validation: What Each Proves}

Synthetic experiments (h-e1, h-m-integrated) prove the gradient abnormality pipeline can differentiate patterns when differences exist. Statistical tests correctly identify divergences, effect sizes compute accurately, and the causal mechanism (background swap $\rightarrow$ GAIA-Z reduction) operates as expected. This validates code correctness and mechanism plausibility in controlled conditions.

Synthetic validation cannot prove whether real Waterbirds minority gradients exhibit abnormality. Synthetic data imposes known GAIA-Z properties (60\% vs 30\% near-zero rates); real gradients may not follow this pattern. If real validation fails (divergence <0.1), synthetic results show the pipeline itself works --- failure would indicate hypothesis falsity, not implementation error.

Proof-of-concept validation (h-m-mitigate MNIST) demonstrates spatial regularization methodology works on a toy dataset with simple spurious feature (color). MNIST overperformance (+23pp) likely due to feature simplicity compared to complex Waterbirds backgrounds. Full 5-seed experiment and real Waterbirds validation needed to confirm robustness and generalization.

Epistemic distinction: Synthetic proves ''The methodology correctly processes gradients when differences exist'' (HIGH confidence). Real would prove ''Real minority gradients exhibit abnormality'' (LOW confidence, untested). This positions our contribution as Tier 3 (methodology validated) with clear path to Tier 2 (pending real empirical validation).

\subsubsection{6.2 GPU Incompatibility as Blocking Factor}

PyTorch 2.0 supports CUDA compute capabilities sm\_37 through sm\_86. NVIDIA H100 NVL GPUs use sm\_90 architecture (Hopper generation). Training aborted with RuntimeError: CUDA error: no kernel image available for sm\_90.

Resolution paths: (1) PyTorch upgrade to 2.1+ (adds sm\_90 support, requires dependency compatibility testing, estimated 1-2 days), or (2) CPU fallback (all experiments runnable on CPU with extended wall-clock time: h-e1 detection ~30h, h-m-integrated mechanism ~50h, h-m-mitigate mitigation ~40h, total ~70-90 GPU-hours equivalent CPU time, 3-4 days wall-clock).

Synthetic validation prioritized during Phase 4 to maximize iteration speed while proving methodology correctness. Real validation deferred to post-submission infrastructure upgrade. Cannot make empirical claims about real Waterbirds worst-group accuracy improvement or minority gradient abnormality without real training. Methodology contribution (pipeline design, statistical framework, causal mechanism) validated independently via synthetic/proof-of-concept tests.

\subsubsection{6.3 MNIST Effectiveness: Interpretation and Prediction}

Observed: +23pp worst-group accuracy improvement (78\% vs 55\%) on MNIST+Color smoke test. Expected: $\geq 10pp$ improvement (gate criterion). Discrepancy: 2.$3\times$ larger than minimum threshold.

Hypothesis: Simple spurious feature (color) vs complex spurious (background texture). MNIST: Binary color channel (red/blue), single-pixel spurious feature --- spatial masking trivial (entire background uniform color). Waterbirds: High-dimensional background (water waves, land textures), multi-region spurious --- spatial masking ambiguous (shoreline pixels part spurious, part informative?).

Gradient regularization penalizes variance in spurious regions. Low-dimensional spurious (color) creates clean binary mask $\rightarrow$ variance penalty effective. High-dimensional spurious (texture) creates noisy mask $\rightarrow$ variance penalty may suppress informative gradients near boundaries.

Predicted Waterbirds performance: Worst-group accuracy improvement 5-15pp (conservative), not 23pp. Spatial regularization expected to outperform ERM (~70\% baseline $\rightarrow$ 75-85\% with regularization), potentially match GroupDRO (80-85\%) but unlikely to beat SCER (~90\%, if reproducible).

\subsubsection{6.4 Contribution Positioning and Tier Assessment}

Current tier: Tier 3 (Methodology Validated).

Validated contributions: (1) Detection methodology: Gradient abnormality pipeline with proven code correctness (synthetic validation). (2) Causal mechanism: 4-step chain validated at synthetic level (correlation test, augmentation test, accuracy check). (3) Mitigation methodology: Spatial gradient regularization framework with proof-of-concept effectiveness (MNIST). (4) Unified framework: Same abnormality mechanism for detection and mitigation (complementary to embedding-based SCER).

Pending for Tier 2 (Empirical Validation): Real Waterbirds detection results (GAIA divergence on real gradients), Real Waterbirds mitigation results (worst-group accuracy comparison to GroupDRO/JTT baselines), Multi-dataset generalization (CelebA, UrbanCars).

Tier 1 unlikely (SCER code unavailable): Cannot claim to beat current SOTA (~90\% worst-group accuracy) without reproducible baseline. Positioning as gradient-based alternative (Tier 2) vs SOTA claim (Tier 1).

Venue recommendation: Workshop (current state): NeurIPS Workshop on Robustness, ICLR Workshop on Spurious Correlations. Framing: ''Methodology proposal with synthetic/proof-of-concept validation, real experiments in progress.'' Main conference (after real validation): NeurIPS, ICML, ICLR. Framing: ''Competitive alternative to GroupDRO with annotation-free detection.''

\subsubsection{6.5 Limitations and Boundary Conditions}

\textbf{L1: Synthetic Validation Only (h-e1, h-m-integrated).} Impact: Code correctness proven, empirical claim (real minority gradients abnormal) unproven. Generalization boundary: Methodology works in controlled setting, real-world applicability pending. Addressability: Resolvable via PyTorch upgrade or CPU training (~70-90h).

\textbf{L2: Toy-Only Mitigation (h-m-mitigate).} Impact: MNIST proof-of-concept effective, Waterbirds generalization unknown. Generalization boundary: Simple spurious (color) vs complex spurious (background) untested. Addressability: Full Waterbirds experiment (~40 GPU hours).

\textbf{L3: No SOTA Comparison.} Impact: Cannot claim superiority over SCER (~90\% worst-group accuracy reported). Generalization boundary: Detection validated, mitigation competitive positioning unproven. Addressability: Independent SCER reproduction pending code release, OR position as complementary gradient-based approach (orthogonal to embedding methods).

\textbf{L4: Single-Seed Proof-of-Concept (h-m-mitigate).} Impact: No statistical significance testing, smoke test may be outlier. Generalization boundary: Methodology works (single run), robustness across seeds unproven. Addressability: 5-seed full experiment with bootstrap confidence intervals.

\textbf{L5: Computational Overhead.} GradCAM extraction adds ~0.5s/sample on CPU. Limits real-time detection (batch inference only). Mitigation: Subsample extraction (32 samples/batch max).

\textbf{L6: Hyperparameter Sensitivity.} Spatial regularization performance depends on $\lambda _init$ and percentile threshold. Grid search required ($3\times 3$ configs = 9 runs). No universal defaults --- dataset-dependent tuning.

\textbf{L7: GradCAM Assumption (A1).} Requires minority samples correctly classified $\geq 60$\% for spatial masking validity. Synthetic validation uses controlled 68\% accuracy; real Waterbirds may fall below threshold (unknown).

\subsubsection{6.6 Relationship to Prior Work}

\textbf{vs GroupDRO:} GroupDRO requires annotations, achieves 80-85\% worst-group accuracy. We propose annotation-free detection + mitigation. Complementary: Our detection could identify groups for GroupDRO input. Competitive: If our mitigation matches GroupDRO worst-group accuracy without annotations (pending validation).

\textbf{vs SCER:} SCER operates on embedding space post-hoc (~90\% worst-group accuracy estimated). We operate on gradient space during training. Complementary: Different intervention points (embeddings vs gradients) may address different failure modes. Competitive: If gradient regularization outperforms embedding refinement (unlikely --- SCER claims SOTA).

\textbf{vs GAIA:} GAIA detects OOD via gradient abnormality. We extend to subpopulation shift (minority groups within in-distribution data). Extension, not competition: Broadens GAIA applicability from distribution shift to spurious correlation detection.

\textbf{vs Adebayo et al.:} Showed attribution fails for unknown spurious. We use abnormality (process), not attribution (result). Orthogonal: Different paradigms --- we bypass attribution ineffectiveness via process-level detection.

Positioning statement: Gradient-based alternative to embedding methods (SCER), complementing annotation-based approaches (GroupDRO). Extends GAIA framework to new domain (spurious correlation detection). Addresses Adebayo's attribution limitations via abnormality paradigm.

\subsubsection{6.7 Future Work}

\textbf{Immediate Next Steps (Address Current Limitations):}

FW1: Real Waterbirds Validation (Detection + Mitigation). Motivation: Synthetic validation proves methodology; real empirical claim pending. Required: PyTorch 2.1+ OR CPU training (~70-90h). Expected outcome: If divergence $\geq 0.2 \rightarrow$ detection validated. If <0.1 $\rightarrow$ hypothesis false, abandon approach. Impact: Tier 3 $\rightarrow$ Tier 2 (empirical validation).

FW2: Full MNIST Experiment (5 seeds $\times$ 50 epochs). Motivation: Smoke test promising but statistically insufficient. Required: ~4 GPU hours. Expected outcome: Confirm worst-group accuracy improvement $\geq 10$\% with bootstrap confidence intervals. Impact: Proof-of-concept $\rightarrow$ statistically robust mitigation methodology.

FW3: GroupDRO Baseline Comparison. Motivation: Competitive positioning requires SOTA comparison. Required: Waterbirds full experiment (3 methods $\times$ 5 seeds $\times$ 300 epochs, ~40 GPU hours). Expected outcome: If worst-group accuracy $\geq GroupDRO \rightarrow$ Tier 2 competitive. If <GroupDRO $\rightarrow$ detection-only contribution.

\textbf{Methodology Extensions:}

FW4: Multi-Dataset Generalization. Test CelebA (gender-hair), UrbanCars (co-occurrence), Medical (demographic bias). Research question: Does gradient abnormality generalize to non-background spurious? Hypothesis: GAIA-Z divergence $\geq 0.15$ on CelebA.

FW5: Joint Gradient + Embedding Regularization. Combine L\_spatial (gradient space) + L\_SCER (embedding space). Research question: Does joint regularization outperform individual methods? Hypothesis: Worst-group accuracy $\geq$ max(SCER, Ours) + 3\%.

FW6: Automatic Percentile Selection. Replace fixed 75th percentile with validation-worst-group-accuracy-based optimization. Research question: Can we eliminate hyperparameter tuning? Hypothesis: Auto-selected percentile achieves worst-group accuracy within 2\% of oracle (grid search optimum).
